\section{Обзор литературы}\label{sec:related_works}
Исследования маршрутизации рассматривают, какую часть вычислений
достаточно выполнить для конкретного запроса и когда дополнительный
эксперт способен улучшить ответ. Для нашей задачи существенны три
группы подходов: правила на основе уверенности первой модели,
сравнительное оценивание полезности перехода и исследования зависимости
обучения/выбора оценивателя от данных. Такое разделение позволяет
сопоставить используемую информацию, обучающие цели и ограничения.

\subsection{Уверенность и сравнительная полезность}
Простые confidence-правила не требуют отдельного обучения роутера,
но оценивают прежде всего надёжность первой модели. В работе
\citet{jitkrittum2023confidence} показано, когда этого достаточно,
и разобраны случаи специализации второго эксперта, шума и сдвига
распределения. Там же используется Diff-01 --- регрессия разности
правильности двух моделей. Поэтому преимущество сравнительного
подхода связано не с самим наличием второго эксперта, а с тем,
содержат ли доступные признаки информацию о неоднородности его пользы.

Для генеративных моделей уверенность дополнительно зависит от способа
агрегации токенных вероятностей и длины ответа.
\citet{gupta2024cascades} показывают ограничения суммы/среднего
и исследуют квантили и embeddings. Для accuracy-задач используется
rescue-only цель, а для перевода --- регрессия разности BLEURT.
Эти цели и информационные режимы нельзя отождествлять. Добавление
признаков второй модели также меняет стоимость решения.

Post-hoc обучение передачи решения учитывает ошибки и стоимость
несовершенного эксперта \citep{narasimhan2022posthoc}.
В AutoMix ожидаемая полезность продолжения оценивается по шумному
сигналу самопроверки в последовательной модели принятия решений
\citep{aggarwal2024automix}. Это более общий механизм каскада,
но самопроверка требует учёта дополнительных затрат.
В unified routing/cascading изучаются условия полезности разных
схем в зависимости от качества оценок \citep{dekoninck2025unified}.
Часть анализа использует искусственный шум истинных quality signals;
это отличается от изменения объёма и состава реального обучения
конкретного оценивателя.

\subsection{Обучающие данные и оценивание разности}
В RouteLLM сравнены роутеры на preference data и её дополнениях
из gold- и judge-разметки \citep{ong2025routellm}. Результаты
демонстрируют роль соответствия источника обучения задаче.
Работа Routing with Generated Data сравнивает query-answer
и query-only методы при разном качестве генератора и фильтрует
запросы по согласию сильных и расхождению слабых моделей
\citep{niu2026routing}. Таким образом, целевой отбор информативных
примеров уже является предметом исследования. В этих работах
выбирается эксперт до получения первого ответа; наш базовый контроль
использует post-answer признаки фиксированной первой модели.

В uplift regression сравниваются прямая оценка эффекта и разность
отдельных моделей \citep{rudas2018linear}, а исследования meta-learners
связывают предпочтительность метода со структурой эффекта и объёмом
данных \citep{kunzel2019metalearners}. Обычно на объекте известен
только один потенциальный исход. При нашей парной разметке оба
ответа наблюдаются, поэтому аргументы об эффективности unpaired
оценивателей нельзя переносить буквально. Например, на одном $X$
с одинаковой ridge-регуляризацией регрессия $S-W$ тождественна
разности двух ridge-прогнозов; различия с logistic-головами возникают
из-за класса функций, потерь и регуляризации.

Особенно близка работа \citet{nyberg2023exploring}: сравниваются
sampling schemes, оцениватели и размеры данных, разбираются
искажение ранжирования и коррекция вероятностей.
Несмотря на пользу большего числа редких событий, изменение состава
train может менять саму оцениваемую условную величину. Это объяснение
необходимо проверять отдельно от статистической поддержки rescue/harm.

\begin{table}[htbp]\centering\small
\caption{Подходы и ограничения, существенные для выбора способа
обучения каскадного маршрутизатора.}\label{tab:related_works}
\begin{tabular}{p{.24\linewidth}p{.31\linewidth}p{.36\linewidth}}
\toprule
Подход & Сильная сторона & Ограничение для нашей задачи \\
\midrule
Confidence \citep{jitkrittum2023confidence} &
Простое правило без отдельной fit-процедуры &
Не моделирует неоднородность пользы второго эксперта \\
Post-hoc gain \citep{jitkrittum2023confidence,gupta2024cascades} &
Использует сравнительную разметку и доступные признаки &
Зависит от объёма данных, представления и выбранной цели \\
Последовательная utility \citep{aggarwal2024automix,dekoninck2025unified} &
Явно сопоставляет качество действий и затраты &
Надёжность quality/verifier signal и его цена влияют на результат \\
Данные для routing \citep{ong2025routellm,niu2026routing} &
Проверяет источники supervision и информативный отбор &
Single-shot режим и разметка отличаются от парного post-answer каскада \\
Sampling effect estimation \citep{nyberg2023exploring} &
Разделяет дисбаланс, искажение разности и коррекцию &
Непарные исходы; прямой отбор индивидуальных rescue/harm недоступен \\
\bottomrule
\end{tabular}
\end{table}

\subsection{Выбор метода и границы текущего исследования}
LLMRouterBench показывает близкие результаты ведущих роутеров
и конкурентоспособность простых решений при едином протоколе
\citep{li2026llmrouterbench}. Эти наблюдения мотивируют проверять,
оправданно ли усложнение, но близкие средние сами по себе не доказывают
эквивалентность. Отдельно следует оценивать качество выбора метода:
такой подход развит для CATE estimator selection
\citep{mahajan2024selection}. В нашей парной постановке качество
политики на dev можно оценить напрямую, без восстановления
ненаблюдаемого контрфактического исхода.

Таким образом, существующая литература задаёт целевую полезность,
сильные базовые методы и известные механизмы чувствительности
к данным. Наше исследование продолжает эту линию контролируемым
сравнением способов обучения при фиксированных экспертах и признаках,
разном составе парных примеров и одинаковой квоте. Требуется
установить, какие различия устойчивы, чем они объясняются и позволяют
ли характеристики train/dev рекомендовать метод на новой задаче.
Таблицы спасений/вреда используются для анализа механизмов,
а не как самостоятельное доказательство новизны или превосходства.
